\begin{proposition}[Proved under stated hypotheses: finite exceptional wild images]
\label{D:finite-exceptional-wild}
Let $X$ be a normal proper integral variety over an algebraically
closed field, let $Y\subsetneq X$ be closed, and put
$U=X\setminus Y$ and $Q=Q(X,Y)$.
Suppose that some normal projective modification $q:X'\to X$
has the following property: for every prime divisor $E$ in
$X'\setminus q^{-1}(U)$, the image $I_E$ of its absolute wild
inertia group in $Q$ is finite.  Then $Q$ is topologically
finitely generated.
It suffices to check the exceptional divisors among these $E$,
because the wild images at the strict transforms of the original
codimension-one boundary are already trivial in $Q$.
\end{proposition}
\begin{proof}
There are finitely many such divisors.  Choose an open normal
subgroup $H\subset Q$ with $H\cap I_E=1$ for all of them.
This is possible by excluding each of the finitely many nonidentity
elements in their finite union.  Normality of $H$ makes the same
condition hold for every conjugate of an $I_E$.
Let $V\to U$ be the connected finite Galois cover classified by
$Q/H$, and let $L=k(V)$.  Normalize $X'$ in $L$, obtaining a
normal projective variety $T$, finite over $X'$.  Over
$U'=q^{-1}(U)$ this normalization is
$V'=V\times_U U'$, which is connected and normal: its generic
fibre is $\operatorname{Spec}L$, and it is finite \'etale over
the normal integral $U'$.
The morphism $V'\to V$ is proper birational.  Consequently
$\pi_1(V')\to\pi_1(V)$ is surjective, since a connected finite
\'etale cover of the normal integral $V$ stays connected after
this birational base change.

We recall the valuation-group comparison needed here.  If $w/v$
is an extension of discrete valuations through a finite separable
field extension $L/K$, choose compatible geometric valuations.
Inside the absolute decomposition group of $v$ we have
\[
 I_w=I_v\cap G_L,\qquad P_w=P_v\cap G_L,
\]
where $I$ and $P$ denote inertia and wild inertia.  The first
identity follows by restricting the action on residue algebraic
closures.  For the second, $P_v\cap G_L$ is a normal pro-$p$
subgroup of $I_w$, and $I_w/(P_v\cap G_L)$ is a subgroup of the
prime-to-$p$ tame inertia quotient $I_v/P_v$; hence it is exactly
the wild inertia subgroup.  This description includes imperfect
residue fields; see the exact sequence in the Stacks Project,
Lemma~58.13.5, tag 0BUA \cite{StacksDTame}.  In characteristic zero all wild groups
are trivial.

Every prime divisor $F$ of $T\setminus V'$ lies above a prime
divisor $E$ of $X'\setminus U'$, because $T\to X'$ is finite.
Under the natural map from $G_L$ to $H$, the preceding identity
puts the image of $P_F$ in $H\cap gI_Eg^{-1}=1$ for a suitable
conjugate.  Thus every finite cover of $V$ classified by a finite
$H$-set, after pullback to $V'$, is tame at every divisorial
boundary component of $T$.

The surjection $\pi_1(V')\to\pi_1(V)\to H$ therefore kills the
closed normal subgroup generated by those wild inertia groups.
It factors as a surjection
\[
 Q(T,T\setminus V')\longrightarrow H.
\]
The source is finitely generated by the projective case of
Proposition~\ref{D:fixed-boundary-positive}.  Hence $H$ is
finitely generated, and so is its finite extension $Q$.

Finally, a proper birational morphism between normal varieties is
an isomorphism at the generic point of every original prime divisor,
so the strict transform has the same discrete valuation.
Its wild inertia image is zero by the definition of $Q(X,Y)$.
This proves the final assertion.
\end{proof}

\begin{remark}
This criterion does not assume that $q$ is an isomorphism over $U$,
nor that the punctured strict-local fundamental groups are finite.
The proof uses only the wild images of the finitely many boundary
valuations on the chosen projective modification.
In any remaining counterexample to finite generation, every such
projective modification must therefore introduce a boundary divisor
with infinite wild image in $Q$.
The criterion by itself is a sufficient-condition refinement, not
an additional separated geometric example.  The elliptic example
of Corollary~\ref{D:schroer-elliptic} supplies the concrete case
where the new images are in fact trivial.
\end{remark}
